\section{CÁC PHƯƠNG PHÁP XỬ LÝ ẢNH}
\subsection{Nguyên tắc chung của khối tăng cường}
Nhiệm vụ của khối tăng cường là biến đổi ảnh thiếu sáng $Y$ thành ảnh $\widehat{X}$ có độ sáng và cấu trúc phù hợp hơn cho quan sát cũng như nhận dạng. Mọi phương pháp đều tuân thủ cùng giao diện: nhận một ảnh RGB $64\times64$ trong $[0,1]$, trả lại ảnh cùng kích thước và không truy cập ảnh tham chiếu. Điều này cho phép đánh giá phương pháp truyền thống và học sâu trong cùng quy trình.

\textbf{Identity} giữ nguyên ảnh tối: $E(Y)=Y$. Đây là mốc bắt buộc để xác định lợi ích thực của tăng cường. Một thuật toán có ảnh đầu ra sáng hơn nhưng Accuracy hoặc Macro-F1 thấp hơn Identity được xem là làm giảm hiệu quả nhận dạng trong điều kiện đó.

\textbf{Clean} là ảnh gốc trước suy giảm. Clean không phải một thuật toán phục hồi và không tham gia lựa chọn bộ tăng cường. Vai trò của nó là tạo mức hiệu năng tham chiếu, cho biết khoảng cách còn lại giữa ảnh sau xử lý với điều kiện đầu vào ban đầu.

\subsection{Hiệu chỉnh Gamma}
\subsubsection{Cơ sở toán học}
Hiệu chỉnh Gamma tăng sáng bằng phép lũy thừa:
\begin{equation}
\widehat{X}=Y^{\alpha},\qquad 0<\alpha<1.
\label{eq:gamma}
\end{equation}
Với cường độ nằm giữa 0 và 1, số mũ nhỏ hơn 1 làm đầu ra lớn hơn đầu vào. Khi $\alpha$ giảm, mức nâng sáng mạnh hơn. Ví dụ, pixel 0,1 trở thành 0,158 với $\alpha=0,8$; 0,251 với $\alpha=0,6$ và 0,398 với $\alpha=0,4$.

\fig{gamma_curves.pdf}{1}{Đường cong Gamma làm tối và tăng sáng trên miền $[0,1]$. Đây là đồ thị giải tích của phép biến đổi, không phải kết quả đo trên tập ảnh.}{fig:gamma}
Trong trường hợp lý tưởng chỉ có Gamma làm tối, lựa chọn $\alpha=1/\gamma_d$ sẽ khôi phục cường độ ban đầu vì $(X^{\gamma_d})^{1/\gamma_d}=X$. Khi có Poisson, Gaussian và bước chặn, phép biến đổi này không còn là nghịch đảo chính xác. Đặc biệt, clipping đã loại bỏ một phần thông tin nên không thể đảo ngược chỉ bằng một hàm điểm.

\subsubsection{Ảnh hưởng đến nhiễu và cách cài đặt}
Độ nhạy cục bộ của hàm tăng sáng đối với đầu vào được mô tả bởi:
\begin{equation}
\frac{\mathrm{d}Y^\alpha}{\mathrm{d}Y}=\alpha Y^{\alpha-1}.
\end{equation}
Khi $Y$ nhỏ và $\alpha<1$, hệ số này có thể lớn. Với nhiễu nhỏ, cách xấp xỉ tuyến tính cho thấy nhiễu đầu vào có thể bị khuếch đại trong vùng tối. Đây là cơ sở để giải thích tại sao giảm Gamma quá mạnh làm hình ảnh sáng rõ nhưng cũng khiến nhiễu hạt nổi bật hơn.

Trong mã nguồn, Gamma được thực hiện trực tiếp bằng phép toán lũy thừa trên mảng RGB số thực, không chuyển qua ảnh 8 bit. Không có phép lọc không gian hoặc khử nhiễu kèm theo. Lưới validation khảo sát $\alpha\in\{0,4;0,6;0,8\}$. Notebook ResNet-18 chọn $\alpha=0,8$, tương ứng mức tăng sáng tương đối nhẹ.

\textbf{Ưu điểm:} cài đặt đơn giản, chi phí thấp, giữ quan hệ đơn điệu của cường độ. \textbf{Hạn chế:} sử dụng cùng quy tắc ở mọi vị trí, không thích ứng với ánh sáng không đồng đều và không phân biệt nét thật với nhiễu. Do áp dụng phi tuyến theo từng kênh, phương pháp cũng không bảo đảm giữ nguyên tỷ lệ giữa các thành phần màu.

\subsection{Cân bằng histogram toàn cục (HE)}
\subsubsection{Nguyên lý phân phối lại mức xám}
Histogram mô tả số pixel xuất hiện tại mỗi mức cường độ. Với ảnh có $L$ mức xám, ký hiệu $n_k$ là số pixel tại mức $k$ và $N$ là tổng số pixel. Histogram chuẩn hóa và hàm phân phối tích lũy được viết:
\begin{equation}
p(k)=\frac{n_k}{N},\qquad C(k)=\sum_{j=0}^{k}p(j).
\end{equation}
Một phép ánh xạ cân bằng điển hình đưa mức xám về:
\begin{equation}
T(k)=\operatorname{round}\!\left((L-1)\frac{C(k)-C_{\min}}{1-C_{\min}}\right),
\label{eq:he}
\end{equation}
trong đó $C_{\min}$ là giá trị CDF khác 0 đầu tiên. Trường hợp ảnh hằng cần được xử lý riêng để tránh mẫu số bằng 0. HE giúp mở rộng khoảng mức xám đang bị tập trung, song không bảo đảm histogram rời rạc sau biến đổi sẽ phẳng hoàn toàn.

\subsubsection{Xử lý trên kênh độ sáng}
Nếu cân bằng độc lập R, G và B, ba phép ánh xạ có thể khác nhau và làm biến đổi mạnh màu sắc. Dự án chuyển ảnh từ RGB sang YCrCb, chỉ áp dụng \texttt{equalizeHist} trên kênh Y, giữ hai kênh màu rồi chuyển lại RGB. Trước khi đổi không gian màu, ảnh được nhân 255, làm tròn và chuyển sang \texttt{uint8}; sau xử lý được chia 255 để trở lại $[0,1]$.

Việc chỉ tác động lên Y giúp hạn chế sai lệch màu so với cân bằng độc lập từng kênh. Tuy nhiên, quá trình lượng tử hóa và chặn miền khi chuyển đổi vẫn có thể làm thay đổi đầu ra. Do đó, không nên mô tả HE như một phép biến đổi bảo toàn màu tuyệt đối.

HE sử dụng thống kê toàn ảnh, nên vùng nền sáng và vùng biển báo tối cùng đóng góp vào một hàm ánh xạ. Trong ảnh có chiếu sáng không đồng đều, cùng một mức xám thuộc các vùng khác nhau vẫn được xử lý giống nhau. Đây là lý do tăng tương phản toàn cục có thể làm một số vùng dễ quan sát hơn nhưng đồng thời bão hòa các vùng vốn đã sáng~\cite{opencv}.

\textbf{Ưu điểm:} không phải chọn tham số tăng sáng, dễ triển khai và nhanh. \textbf{Hạn chế:} có thể tăng nhiễu, thay đổi phân bố sáng quá mạnh và làm mất thông tin trong vùng bị bão hòa. Với crop nhỏ, histogram cũng nhạy với tỷ lệ nền và kích thước phần biển báo trong ảnh.

\subsection{Cân bằng histogram thích nghi có giới hạn tương phản (CLAHE)}
\subsubsection{Từ xử lý toàn cục đến xử lý cục bộ}
CLAHE chia ảnh thành một lưới ô và xây dựng phép ánh xạ histogram cho từng ô. Cách tiếp cận cục bộ giúp xử lý các vùng có điều kiện chiếu sáng khác nhau. Tuy nhiên, nếu cân bằng cục bộ không giới hạn, một vùng gần như đồng nhất có thể bị kéo giãn tương phản mạnh và làm nhiễu nổi rõ. CLAHE giải quyết điều này bằng cách giới hạn chiều cao histogram trước khi xây dựng CDF.

Quy trình gồm bốn bước: tính histogram mỗi ô; cắt các bin vượt ngưỡng; phân phối lại lượng dư; xây dựng ánh xạ và nội suy giữa các ô. Nội suy ở ranh giới giúp giảm hiện tượng xuất hiện các khối vuông sau xử lý. Trong triển khai của dự án, toàn bộ phép này được áp dụng trên kênh Y, tương tự HE.

\subsubsection{Ý nghĩa tham số và ảnh hưởng kích thước ảnh}
Hai tham số chính là \texttt{clipLimit} và \texttt{tileGridSize}. Tham số thứ nhất điều khiển mức giới hạn histogram theo quy ước OpenCV; không phải ngưỡng cường độ sáng của pixel. Tham số thứ hai là số ô theo hai chiều, không phải kích thước mỗi ô tính bằng pixel. Với ảnh $64\times64$, lưới $8\times8$ tạo ô khoảng $8\times8$ pixel, còn lưới $4\times4$ tạo ô khoảng $16\times16$ pixel.

Ở ô quá nhỏ, số mẫu dùng để ước lượng histogram ít, dễ khiến thống kê bị ảnh hưởng bởi nhiễu và từng nét của biển báo. Lưới thưa hơn có thể ổn định hơn nhưng mức thích ứng không gian cũng giảm. Do vậy, cấu hình tốt trên ảnh độ phân giải cao chưa chắc tốt khi chuyển sang crop $64\times64$.

Lưới validation khảo sát ba giá trị clip limit $\{1,2,4\}$ kết hợp hai lưới $\{4\times4,8\times8\}$. Notebook chọn clip limit bằng 1 và lưới $4\times4$. Đây là kết quả lựa chọn bằng khả năng nhận dạng trên validation, không phải tham số được chọn bằng cách nhìn ảnh test.

\fig{parameters.png}{0.88}{Ảnh hưởng của Gamma và clip limit trên cùng ảnh tối mức vừa. Hàng dưới giữ lưới CLAHE $8\times8$ để minh họa riêng tham số clip limit; đây không phải cấu hình CLAHE được chọn cuối cùng.}{fig:parameters}
Với ô nhỏ, các giá trị clip limit danh nghĩa khác nhau có thể tạo cùng ngưỡng hiệu dụng. Với ô $8\times8$ có 64 pixel và 256 bin, việc quy đổi ngưỡng cắt về số nguyên trong OpenCV có thể đưa các giá trị này về cùng ngưỡng hiệu dụng. Vì vậy, đổi tham số danh nghĩa chưa chắc tạo ra đầu ra khác trên ảnh nhỏ. Nhận định này phù hợp với cài đặt thư viện và quan sát thí nghiệm, không nên suy rộng thành tính chất chung cho mọi kích thước ảnh.

\textbf{Ưu điểm:} tăng tương phản theo vùng, kiểm soát tốt hơn sự khuếch đại histogram và có chi phí tính toán thấp hơn nhiều so với các bộ phục hồi sâu. \textbf{Hạn chế:} vẫn có thể làm nhiễu rõ hơn; chất lượng phụ thuộc tương tác giữa clip limit, lưới và độ phân giải.

\subsection{Retinex đa tỉ lệ (Multi-Scale Retinex)}
\subsubsection{Mô hình phản xạ và chiếu sáng}
Retinex xem ảnh quan sát là tích của thành phần phản xạ $R$ và chiếu sáng $L$:
\begin{equation}
I(x,y)=R(x,y)L(x,y).
\end{equation}
Nếu có thể ước lượng thành phần chiếu sáng biến đổi chậm, ta có thể giảm ảnh hưởng của nó và làm rõ cấu trúc phản xạ. Trong miền log, tích trở thành tổng; một dạng Single-Scale Retinex sử dụng ảnh làm mờ Gaussian để xấp xỉ thành phần chiếu sáng:
\begin{equation}
r_s=\log(Y+\epsilon)-\log(G_{\sigma_s}*Y+\epsilon).
\end{equation}
Với nhiều độ lệch chuẩn Gaussian, kết quả trung bình là:
\begin{equation}
r=\frac{1}{S}\sum_{s=1}^{S}r_s,\qquad \epsilon=10^{-6}.
\label{eq:msr}
\end{equation}
Các độ lệch chuẩn nhỏ nhấn mạnh biến đổi cục bộ, trong khi độ lệch chuẩn lớn thể hiện cấu trúc chiếu sáng rộng hơn. Kết hợp nhiều tỉ lệ giúp giảm phụ thuộc vào một mức làm mờ duy nhất.

\subsubsection{Chuẩn hóa đầu ra trong dự án}
Giá trị $r$ trong miền log chưa nằm trong $[0,1]$. Mã nguồn lấy percentile 1\% và 99\% trên toàn mảng kết quả, ký hiệu là $q_1,q_{99}$, rồi chuẩn hóa:
\begin{equation}
\widehat{X}=\operatorname{clip}\!\left(\frac{r-q_1}{q_{99}-q_1},0,1\right).
\label{eq:retinexnorm}
\end{equation}
Nếu $q_{99}-q_1<10^{-6}$, hàm trả lại ảnh đầu vào để tránh chia cho khoảng quá nhỏ. Hai cấu hình tỉ lệ được thử là $[3,10,20]$ và $[5,15,30]$; notebook chọn cấu hình thứ hai.

Đây là MSR giản lược, không có bước phục hồi màu riêng của MSRCR và không có khối khử nhiễu chuyên biệt. Việc chuẩn hóa percentile làm kết quả phụ thuộc vào phân bố của chính ảnh đang xử lý. Một số pixel bị kéo đến 0 hoặc 1; quan hệ màu và độ sáng tuyệt đối có thể thay đổi đáng kể so với ảnh gốc.

\subsubsection{Vì sao Retinex nhạy với nhiễu?}
Vùng tối có giá trị $Y$ nhỏ, trong khi đạo hàm của hàm log là $1/(Y+\epsilon)$. Các dao động nhỏ của nhiễu vì vậy có thể tạo chênh lệch lớn trong miền log. Phép trừ ảnh làm mờ còn nhấn mạnh thành phần biến thiên nhanh, bao gồm cả biên thật lẫn nhiễu. Sau bước kéo giãn percentile, những dao động này có thể tiếp tục được khuếch đại.

Khi ảnh chỉ bị Gamma làm tối, MSR có thể làm rõ cấu trúc tốt. Khi thêm nhiễu, cùng cấu hình có thể tạo màu giả và hạt rõ trên nền. Thí nghiệm ablation ở phần sau kiểm tra trực tiếp giả thuyết này bằng cách giữ nguyên $\gamma_d=2,2$ và thay đổi thành phần nhiễu.

\subsection{Các phương pháp tăng cường dựa trên học sâu}
\subsubsection{Zero-DCE}
Zero-DCE ước lượng các đường cong điều chỉnh cường độ phụ thuộc ảnh. Một dạng đường cong được mô tả trong nghiên cứu gốc là $I'=I+A I(1-I)$, với tham số $A$ được mạng dự đoán theo vị trí và được áp dụng lặp. Phương pháp học cách điều chỉnh ảnh bằng các ràng buộc không cần ảnh tham chiếu ghép cặp~\cite{zerodce}.

Trong dự án, Zero-DCE được dùng ở chế độ suy luận với checkpoint \texttt{Epoch99.pth}. Không tối ưu lại mạng trên GTSRB. Adapter chọn ảnh tăng cường từ đầu ra của mô hình, đưa về miền hợp lệ và cùng kích thước ban đầu. Vì đường cong tác động lên ảnh chứa nhiễu, đầu ra sáng hơn vẫn cần được kiểm tra về hạt và độ rõ nét số.

\subsubsection{Retinexformer}
Retinexformer kết hợp biểu diễn chiếu sáng với Transformer để phục hồi ảnh thiếu sáng. Khung một giai đoạn của phương pháp sử dụng thông tin chiếu sáng nhằm hỗ trợ xử lý các suy giảm và quan hệ giữa các vùng ảnh~\cite{retinexformer}. Báo cáo sử dụng nó như một đối chứng học sâu cho MSR, không đồng nhất hai phương pháp chỉ vì cùng dựa trên ý tưởng Retinex.

Checkpoint được dùng là bản LOL-v1 do tác giả công bố. Adapter nạp kiến trúc và trọng số, chạy suy luận, sau đó lấy phần đầu ra tương ứng ảnh gốc. Ảnh cần được đệm đến bội số của 4 nếu kích thước chưa phù hợp; với crop $64\times64$, điều kiện này đã được thỏa mãn.

\subsubsection{Tăng cường dựa trên diffusion trong miền wavelet}
Phương pháp diffusion được chọn là Wavelet-based Diffusion LLIE. Nghiên cứu gốc kết hợp mô hình diffusion có điều kiện với biểu diễn wavelet và cơ chế phục hồi thành phần tần số cao~\cite{diffusion}. Cách tiếp cận này hướng tới cân bằng thông tin chiếu sáng và chi tiết của ảnh.

Trong triển khai, mạng dùng trọng số trung bình trượt (EMA), 10 bước lấy mẫu và hạt giống suy luận được xác định theo mã ảnh. Nếu cần, ảnh được đệm lặp biên đến bội số của 32; đầu ra lấy từ trường \texttt{pred\_x}, cắt về kích thước ban đầu và chặn $[0,1]$. Nhiều bước suy luận làm chi phí lớn hơn đáng kể so với Gamma hoặc CLAHE.

Ba mô hình sâu đều được dùng nguyên trọng số tác giả. Vì dữ liệu học trước khác miền biển báo GTSRB và khác mô hình nhiễu tổng hợp đang khảo sát, không có cơ sở mặc định rằng chúng sẽ vượt mọi phương pháp truyền thống. Đây là một nội dung cần kiểm chứng bằng kết quả thực nghiệm thay vì suy đoán theo độ phức tạp của mô hình.

\subsection{Minh họa trước và sau xử lý}
\fig{classical_default.png}{1}{Ảnh và histogram sau các phép xử lý truyền thống với tham số mặc định: Gamma 0,6; CLAHE clip limit 2, lưới $8\times8$; MSR $[3,10,20]$. Hình dùng để giải thích cơ chế, không phải kết quả của bộ tham số đã chọn.}{fig:classical}
Hình~\ref{fig:classical} cho thấy cùng một ảnh tối có thể được biến đổi thành nhiều phân bố sáng khác nhau. Gamma nâng phần lớn pixel tối; HE phân phối lại mức sáng trên toàn ảnh; CLAHE phụ thuộc từng vùng; MSR có thể tạo thay đổi mạnh về màu và độ tương phản. Việc histogram trải rộng sau xử lý chỉ chứng minh phép ánh xạ đã thay đổi phân bố, chưa chứng minh nội dung có ích đã được phục hồi.

\fig{enhancers_selected.png}{0.82}{Minh họa bảy phương pháp trên một crop: Gamma 0,8; CLAHE clip 2, lưới $4\times4$; MSR $[5,15,30]$. Clean là tham chiếu, Identity là ảnh tối. Hình giải thích biến đổi ảnh, không phải bằng chứng phân loại của bộ tham số cuối.}{fig:all}
Trong Hình~\ref{fig:all}, ảnh Zero-DCE và Retinexformer sáng lên rõ, nhưng nhiễu và sự biến đổi màu cũng cần được quan sát. MSR làm nền thay đổi mạnh, còn diffusion không nhất thiết tạo ảnh sáng nhất. Những nhận xét này chỉ có giá trị minh họa trên một mẫu; xếp hạng thuật toán phải dựa trên số liệu tổng hợp và hiệu quả nhận dạng ở các phần tiếp theo.

\subsection{Đánh giá chất lượng bằng PSNR và SSIM}
\subsubsection{Sai số bình phương trung bình và PSNR}
Với $X$ là ảnh gốc, $\widehat{X}$ là ảnh xử lý và $N=HWC$ là số phần tử RGB:
\begin{equation}
\operatorname{MSE}=\frac{1}{N}\sum_{i=1}^{N}(X_i-\widehat{X}_i)^2,
\qquad \operatorname{PSNR}=10\log_{10}\frac{1}{\operatorname{MSE}}.
\end{equation}
Công thức dùng biên độ cực đại bằng 1 vì ảnh đã được chuẩn hóa. PSNR cao tương ứng sai số pixel nhỏ. MSE được tính với độ chính xác số thực 64 bit trong mã nguồn để giảm sai số cộng dồn. Nếu ảnh trùng hoàn toàn với tham chiếu, MSE bằng 0 và PSNR là vô hạn, không thay bằng một giá trị hữu hạn tùy ý.

PSNR được tính từng ảnh rồi lấy trung bình. Trung bình của log sai số khác log của sai số trung bình, nên không được gộp toàn bộ pixel của tập dữ liệu thành một MSE duy nhất rồi gọi đó là PSNR trung bình. Chỉ số cũng nhạy với thay đổi độ sáng toàn cục, dù ảnh sau thay đổi vẫn có thể giữ được các nét phân lớp.

\subsubsection{Độ tương đồng cấu trúc SSIM}
SSIM xem xét thống kê cục bộ của hai ảnh, gồm độ sáng, độ tương phản và mức tương đồng cấu trúc~\cite{ssim}. Dạng thường dùng là:
\begin{equation}
\operatorname{SSIM}(x,y)=
\frac{(2\mu_x\mu_y+C_1)(2\sigma_{xy}+C_2)}
{(\mu_x^2+\mu_y^2+C_1)(\sigma_x^2+\sigma_y^2+C_2)}.
\end{equation}
Trong đó $\mu$ là trung bình cục bộ, $\sigma^2$ là phương sai, $\sigma_{xy}$ là hiệp phương sai và các hằng số $C_1,C_2$ ổn định phép chia. Dự án dùng cửa sổ $7\times7$, miền giá trị 1, xử lý ba kênh màu, $K_1=0,01$, $K_2=0,03$, không dùng trọng số Gaussian và dùng hiệp phương sai mẫu.

SSIM bổ sung thông tin mà PSNR không thể hiện đầy đủ, song vẫn không thay thế metric nhận dạng. Nét chữ số bị sai ở một vùng nhỏ có thể không làm giảm mạnh điểm trung bình toàn ảnh. Khi ảnh đêm thực không có ảnh tham chiếu đồng vị trí, cũng không thể áp dụng trực tiếp PSNR/SSIM theo giao thức hiện tại.
